\begin{textsample}{NEG Dim 3 – 2019\_c5 – Score 2.00 – 2019\_c5/enade\_2019\_med\_cf41be24ef.txt}  \label{ex:f3_neg_191}
CAPÍTULO 5 DISTRIBUIÇÃO DOS CONCEITOS 5.1 PANORAMA NACIONAL DA DISTRIBUIÇÃO DOS CONCEITOS O Conceito Enade é um indicador de qualidade que avalia os cursos por intermédio dos desempenhos dos estudantes a partir dos resultados do Enade.
Os valores possíveis vão de 1 ( pior situação ) a 5 ( melhor situação ).
Os cursos com apenas um ou sem nenhum concluinte participante não obtêm o Conceito Enade, ficando Sem Conceito ( SC ).
Na Tabela 5.1, são apresentadas a quantidade e a distribuição de cursos de Medicina participantes do Enade/2019, por faixa de conceito e Grande Região.
Enfatiza -se, mais uma vez, que a diferença entre os cursos tabulados neste capítulo e no capítulo 2 corresponde aos cursos que não foram avaliados, ou seja, àqueles, em princípio, sem alunos concluintes que tivessem sido inscritos para a prova.
Observando -se os dados da Tabela 5.1, nota -se que, dos 232 cursos participantes, 90 ( 38,8 \% ) classificaram -se com conceito 4, o valor modal.
Esse foi também o conceito modal nas regiões Nordeste ( 45,8 \% ), Sul ( 52,6 \% ) e Centro-Oeste ( 47,6 \% ).
O conceito 3 foi o segundo mais frequente em nível nacional ( 35,8 \%, correspondendo a 83 cursos ), e o conceito 5, o terceiro ( 12,1 \%, correspondendo a 28 cursos ).
Houve, ainda, 18 cursos ( 7,8 \% ) que receberam conceito 2, e 13 cursos ( 5,6 \% ) que receberam conceito 1.
Dos 232 cursos de Medicina, nenhum ficou Sem Conceito ( SC ).
A região Norte participou com 17 cursos ou 7,3 \% do total nacional.
Desses, sete cursos receberam o conceito 3, o que equivale a 41,2 \% do total regional, além de ser o conceito modal na região.
A três cursos ( 17,6 \% ) atribuiu -se conceito 2, e a dois cursos ( 11,8 \% ), conceito 4.
Cinco cursos receberam conceito 1.
Nenhum curso ficou Sem Conceito ( SC ) ou recebeu conceito 5.
A região Nordeste participou com 59 cursos ou 25,4 \% do total nacional.
Nessa região, 27 cursos ( 45,8 \% em termos regionais ) obtiveram conceito 4, o conceito modal para a região.
Os conceitos 2 e 5 foram atribuídos, respectivamente, a três e a oito cursos ( 5,1 \% e 13,6 \%, respectivamente ).
Ainda nessa região, 21 cursos ( 35,6 \% ) foram avaliados com o conceito 3.
Nenhum curso ficou Sem Conceito ( SC ) ou recebeu conceito 1.
Dos 97 cursos participantes da região Sudeste ( 41,8 \% do total nacional ), o conceito 3 foi o conceito modal, tendo sido obtido por 38 cursos ( 39,2 \% ).
O conceito 1 foi atribuído a sete cursos ( 7,2 \% ), e o conceito 2, a nove cursos ( 9,3 \% ).
Receberam os conceitos 4 e 5, respectivamente, 31 cursos ( 32,0 \% ) e 12 cursos ( 12,4 \% ).
Nessa região, nenhum curso ficou Sem Conceito ( SC ).
Os 38 cursos da região Sul corresponderam a 16,4 \% do total nacional.
Como já indicado, houve predominância do conceito 4, atribuído a 52,6 \% dos cursos da região ( conceito modal ), ou seja, atribuído a 20 dos 38 cursos participantes na região Sul.
O conceito 2 foi atribuído a dois cursos ( 5,3 \% ), conceito 3 a 12 cursos ( 31,6 \% ), conceito 5 a quatro cursos ( 10,5 \% ).
Nenhum curso recebeu conceito 1 ou ficou Sem Conceito ( SC ).
Já dos 21 cursos participantes na região Centro-Oeste ( 9,1 \% do total nacional ), dez ( 47,6 \% em termos regionais ) receberam conceito 4, o conceito modal.
Os conceitos 1 e 2 foram atribuídos a um curso cada ( 4,8 \% ).
Cinco cursos receberam conceito 3, conceito 5 foi atribuído a quatro cursos.
Nenhum curso ficou Sem Conhecimento ( SC ).
No Gráfico 5.1, é apresentada a distribuição cumulativa do Conceito Enade segundo Grande Região.
Em um gráfico de distribuição cumulativa, em cada valor de abscissa ( eixo x ), o valor apresentado no eixo das ordenadas ( eixo y ) é o valor acumulado da variável ( conceito Enade ).
Observando a linha vertical pontilhada, que ajuda a visualizar o percentual de conceitos acumulados até o valor 3, ou seja, conceitos 1, 2 e 3, pode -se notar que, por exemplo, a região Norte ( linha azul ) apresenta mais de 88 \% dos cursos com conceito 3 ou abaixo deste valor.
A região Centro-Oeste ( linha laranja ), por outro lado, apresenta menos de 34 \% dos cursos na mesma situação : com conceito 3 ou abaixo.
Destaca -se que na região Norte ( linha azul ) o conceito 5 não foi alcançado por nenhum curso, todos os cursos têm conceito 4 ou abaixo disso e a poligonal cumulativa alcança 100 \% no conceito 4.
Em linhas gerais, em um gráfico de distribuição cumulativa, poligonais mais à esquerda representam Grandes Regiões com pior distribuição, e poligonais mais à direita, Grandes Regiões com melhores conceitos.
A região Norte ( linha azul ) apresenta os cursos com conceitos mais baixos, ao passo que as regiões Centro-Oeste ( linha laranja ) e Sul ( linha rosa ) apresentam as distribuições com valores maiores.
Cursos em IES situadas nas regiões Sudeste ( linha vermelha ) e Nordeste ( linha verde ) apresentam uma situação intermediária. 5.2 CONCEITOS POR CATEGORIA ADMINISTRATIVA E GRANDE REGIÃO Na Tabela 5.2, é apresentada a distribuição dos cursos de Medicina participantes do Enade/2019 por Categoria Administrativa, de acordo com os conceitos por eles alcançados, segundo as Grandes Regiões brasileiras.
Dos 232 cursos participantes, 100 ( 43,1 \% ) eram ministrados em Instituições Públicas, e 132 ( 56,9 \% ), em Instituições Privadas.
De acordo com as informações da Tabela 5.2, em termos nacionais, entre Instituições Públicas e Privadas, observa -se que, dos 28 cursos avaliados com conceito 5, 24 eram oferecidos em IES Públicas, e quatro, em IES Privadas.
Dos 100 cursos participantes de IES Públicas, o conceito 4 foi o valor modal, atribuído a 51 cursos ( 51,0 \% ).
Dos demais cursos de Instituições Públicas participantes, quatro cursos foram avaliados com conceito 2, 19 cursos ( 19,0 \% ) foram avaliados com conceito 3 e, 24 cursos ( 24,0 \% ), com conceito 5.
Dois cursos ficaram com conceito 1, e nenhum curso ficou Sem Conhecimento ( SC ).
Na rede privada, o conceito modal foi 3, atribuído a 64 cursos, o correspondente a 48,5 \% dos 132 cursos da categoria.
Dos demais cursos participantes, 11 ( 8,3 \% ) receberam conceito 1, e 14 ( 10,6 \% ), conceito 2.
O conceito 4 foi atribuído a 39 cursos ( 29,5 \% ), e o conceito 5, a quatro cursos ( 3,0 \% ).
Nesta Categoria Administrativa, nenhum curso ficou Sem Conceito ( SC ).
Na análise por região, observa -se que, na região Norte, as Instituições Privadas participaram com oito cursos ( 47,1 \% do total regional ), dos quais nenhum ficou Sem Conceito ( SC ), e nenhum recebeu os conceitos 4 ou 5.
O conceito modal para as Instituições Privadas na região foi 3, atribuído a quatro cursos, correspondendo a 50,0 \% dessa categoria na região.
Um curso foi avaliado com conceito 2, e três cursos, com conceito 1.
As Instituições Públicas Na região Norte, as Instituições Públicas participaram com nove cursos na região Norte ( 52,9 \% do total regional ), aos quais foram atribuídos os conceitos 1, 2, 3 e 4.
Nenhum curso ficou Sem Conceito ( SC ) ou com conceito 5 nessa combinação de categoria e região.
Na região Nordeste, a rede privada concentrou 24 dos 59 cursos participantes, o equivalente a 40,7 \% do total da região.
Dentre os cursos oferecidos por IES Privadas no Nordeste, 15 cursos foram avaliados com conceito 3, o conceito modal.
Dois cursos obtiveram conceito 2, seis cursos receberam conceito 4 e um curso ficou com conceito 5.
Nessa combinação de categoria e região, nenhum curso ficou Sem Conceito ( SC ) e nenhum recebeu conceito 1.
As Instituições Públicas da região Nordeste participaram com 35 cursos ( 59,3 \% do total da região ).
Desses, 21 obtiveram conceito 4, o valor modal, seguidos por outros sete cursos, que receberam conceito 5.
Seis cursos obtiveram conceito 3, e um curso recebeu conceito 2.
Nenhum curso ficou Sem Conceito ( SC ) ou com o conceito 1.
Na região Sudeste, a proporção de cursos da rede privada, 71,1 \%, correspondeu a 69 dos 97 cursos participantes.
Nessa categoria e região, o conceito modal foi 3, atribuído a 34 cursos, e nenhum ficou Sem Conceito ( SC ).
Os demais foram avaliados com conceito 1 ( sete cursos ), conceito 2 ( nove cursos ), conceito 4 ( 18 cursos ) e conceito 5 ( um curso ).
Dos 28 cursos oferecidos por Instituições Públicas na região Sudeste, o conceito modal foi 4, atribuído a 13 cursos.
Nessa região, 11 cursos receberam conceito 5, e quatro receberam conceito 3.
Nenhum curso recebeu conceitos 1 e 2 ou ficou Sem Conceito ( SC ).
As Instituições Privadas concentraram 23 dos 38 cursos participantes da região Sul, 60,5 \% do total regional.
Desses, dez cursos obtiveram conceito 4, o conceito modal.
Nesta combinação de Categoria Administrativa e Grande Região, dois cursos receberam conceito 2, outros dois, conceito 5, e ainda outros nove, conceito 3.
Nenhum curso recebeu conceito 1 ou ficou Sem Conceito ( SC ).
As Instituições Públicas na região Sul participaram com 15 cursos ( 39,5 \% dos cursos da região ), aos quais foram atribuídos os conceitos 3, 4 e 5.
Nenhum curso da região Sul oferecido por IES Pública ficou Sem Conceito ( SC ) ou recebeu os conceitos 1 e 2.
Na região Centro-Oeste, oito dos 21 cursos participantes eram de Instituições Privadas ( 38,1 \% em termos regionais ).
Desses, cinco receberam conceito 4, o conceito modal.
Dos demais cursos, dois receberam conceito 3, e o outro, conceito 1.
Nenhum curso ficou Sem Conceito ( SC ), e nenhum curso recebeu conceito 2 e 5.
Quanto aos 13 cursos oferecidos por Instituições Públicas na região Centro-Oeste ( 61,9 \% do total regional ), cinco cursos apresentaram conceito 4, conceito modal, e quatro cursos receberam conceito 5.
O conceito 3 foi atribuído a três cursos e o conceito 2 a um curso.
Nenhum curso ficou Sem Conceito ( SC ) ou recebeu conceito 1.
No Gráfico 5.2, apresenta -se a distribuição do Conceito Enade, segundo a Categoria Administrativa da IES.
Os cursos em IES Públicas ( linha azul ) apresentam uma poligonal mais à direita do que os em IES Privadas ( linha verde ) e, consequentemente, uma distribuição de Conceitos Enade com valores maiores. 5.3 CONCEITOS POR ORGANIZAÇÃO ACADÊMICA E GRANDE REGIÃO Na Tabela 5.3, encontra -se a distribuição dos conceitos atribuídos aos cursos participantes do Enade/2019, na Área de Medicina, por Organização Acadêmica, segundo as Grandes Regiões brasileiras.
Dos 232 cursos de Medicina participantes, 155 eram oferecidos em Universidades, 40 em Centros Universitários, 37 em Faculdades e nenhum em CEFET/IF.
Esta distribuição corresponde a, respectivamente, 66,8 \%, 17,2 \%, 15,9 \% e 0,0 \% dos cursos.
De acordo com os dados apresentados, de todos os 28 cursos avaliados com conceito 5, 26 eram vinculados a Universidades.
Os cursos vinculados a esse tipo de Organização Acadêmica tiveram o conceito 4 como conceito modal, atribuído a 67 cursos ( 43,2 \% ).
Os demais cursos vinculados a Universidades receberam os conceitos 1 ( nove cursos ), 2 ( dez cursos ), 3 ( 43 cursos ) e conceito 5 ( 26 cursos, como já mencionado ).
Nenhum curso ficou Sem Conceito ( SC ).
Dos cursos em Centros Universitários, o conceito modal foi 3, atribuído a 21 cursos ( 52,5 \% ).
Dois cursos vinculados a esse tipo de Organização Acadêmica ficaram com conceito 1, cinco com conceito 2, 12 com conceito 4.
Nenhum curso ficou Sem Conceito ( SC ) ou recebeu conceito 5.
Dos 37 cursos mantidos por Faculdades, 19 ( 51,4 \% ) receberam conceito 3, o conceito modal.
O segundo conceito mais frequentemente atribuído aos cursos oferecidos por essa categoria de IES foi o conceito 4, recebido por 11 cursos.
Na sequência, três cursos obtiveram conceito 2, e os conceitos 1 e 5 foram atribuídos a dois cursos cada.
Dos cursos oferecidos em Faculdades, nenhum ficou Sem Conceito ( SC ).
Considerando -se, separadamente, as regiões brasileiras, verifica -se que, na região Norte, as Universidades participaram com dez cursos avaliados.
Os conceitos 1 e 3 foram concedidos a três cursos cada, os conceitos 2 e 4 foram atribuídos a dois cursos cada.
Nenhum curso ficou Sem Conceito ( SC ) ou recebeu conceito 5.
Os Centros Universitários da região Norte foram representados por quatro cursos, que receberam conceitos 3 ( três cursos ) e 2 ( um curso ).
Nenhum curso ficou Sem Conceito ( SC ) ou recebeu os demais conceitos nesta combinação de categoria e região.
As Faculdades As Faculdades participaram com três cursos na região Norte, que receberam os conceitos 1 ( dois cursos ) e 3 ( um curso ).
Nenhum curso ficou Sem Conceito ( SC ) ou recebeu os demais conceitos.
Na região Nordeste, as Universidades participaram com 41 dos 59 cursos da Área de Medicina oferecidos.
Nenhum desses ficou Sem Conceito ( SC ).
O conceito modal 4 foi atribuído 22 cursos.
Os conceitos 2, 3 e 5 foram atribuídos, respectivamente, a dois, nove e oito cursos.
Nenhum curso oferecido por Universidades no Nordeste recebeu conceito 1.
Os Centros Universitários contaram com 11 cursos participantes na região Nordeste, oito dos quais receberam o conceito modal 3.
Para os demais cursos, foi atribuído o conceito 4 ( três cursos ).
Nenhum curso oferecido por Centros Universitários no Nordeste ficou Sem Conceito ( SC ) ou recebeu os conceitos 1, 2 e 5.
As Faculdades foram representadas por sete cursos na região Nordeste, quatro receberam conceito 3, o conceito modal.
O conceito 4 foi atribuído a dois cursos e conceito 2 a um curso.
Nenhum curso ficou Sem Conceito ( SC ) ou recebeu outro conceito.
Na região Sudeste, as Universidades concentraram 55 dos 97 cursos de Medicina da região.
Dos cursos oferecidos em Universidades, nessa região, o conceito modal foi 4, atribuído a 19 cursos.
Os demais cursos receberam os conceitos 1 ( seis cursos ), 2 ( quatro cursos ), 3 ( 16 cursos ) e 5 ( dez cursos ).
Nenhum curso ficou Sem Conceito ( SC ).
Os Centros Universitários participaram com 18 cursos na região Sudeste, dos quais nove obtiveram o conceito modal 3.
O conceito 1 foi atribuído a um curso, o conceito 2, a três cursos, e o conceito 4, a cinco cursos.
Nenhum curso ficou Sem Conceito ( SC ) nessa combinação de organização acadêmica e região.
Já as Faculdades foram representadas por 24 cursos na região Sudeste, nenhum ficou Sem Conceito ( SC ) ou recebeu conceito 1.
O conceito modal 3 foi atribuído a 13 cursos.
Aos demais cursos foram atribuídos os conceitos 2 ( dois cursos ), 4 ( sete cursos ) e 5 ( dois cursos ).
Dos 38 cursos da região Sul, 33 eram mantidos por Universidades, dos quais o conceito modal foi 4, atribuído a 18 cursos.
Os demais cursos receberam os conceitos 3 ( dez cursos ), 5 ( 4 cursos ) e 2 ( um curso ).
Nesse tipo de organização, nenhum curso da região Sul ficou Sem Conceito ( SC ) ou recebeu conceito 1.
Dos três cursos participantes de Centros Universitários da região Sul, nenhum ficou sem Conceito ( SC ), e nenhum recebeu os conceitos 1 ou 5.
Os conceitos restantes, 2, 3 e 4, foram atribuídos a um curso, cada.
Quanto aos dois cursos vinculados a Faculdades na região Sul, um recebeu conceito 3 e outro recebeu conceito 4.
Já os CEFET/IF não participaram na oferta de cursos na região Sul. ( SC ).
Os conceitos 2, 3, 4 e 5 foram atribuídos, respectivamente, a um, cinco, seis e quatro cursos.
Nesse tipo de organização, nenhum curso recebeu conceito 1.
Os Centros Universitários da região Centro-Oeste contaram com quatro cursos, três avaliados com conceito 4, e um avaliado com conceito 1.
Os demais conceitos não foram utilizados.
Nenhum curso ficou Sem Conceito ( SC ).
As Faculdades foram representadas por um curso que recebeu conceito 4.
Já os CEFET/IF não participaram com oferta de cursos na região Centro-Oeste.
No Gráfico 5.3, apresenta -se a distribuição do Conceito Enade dos cursos de Medicina, segundo a Organização Acadêmica da IES.
Os cursos em Universidades ( linha azul ) apresentam uma poligonal mais à direita do que as demais, denotando uma melhor distribuição de Conceitos Enade.
No outro extremo, com distribuição muito semelhante entre si, encontram -se os cursos mantidos por Centros Universitários ( linha verde ) e Faculdades ( linha vermelha ), apresentando as respectivas poligonais mais à esquerda e, consequentemente, distribuição de Conceitos Enade com menores valores.

% matched lemmas: 
\end{textsample}
